%=======================================================================
\chapter{Jeu de données et prétraitement}

\section{Présentation du dataset}

\subsection{Source du dataset}

Dans le cadre de ce projet, nous avons utilisé un jeu de données d'images d'Imagerie par
Résonance Magnétique (IRM) cérébrale destiné à la classification multiclasses des tumeurs
cérébrales. Ce dataset est largement utilisé dans la littérature scientifique pour
l'évaluation des algorithmes de classification médicale et contient des images représentant
différents types de tumeurs cérébrales.

Les images sont issues d'examens IRM réalisés sur des patients présentant diverses
pathologies tumorales. Chaque image est associée à une étiquette indiquant la catégorie de
tumeur correspondante, permettant ainsi l'utilisation de méthodes d'apprentissage supervisé.

\subsection{Composition du dataset}

Afin de garantir un apprentissage équilibré et d'éviter les biais liés au déséquilibre des
classes, un sous-ensemble équilibré de \NTotal{} images a été constitué. Le jeu de données
est réparti équitablement entre trois catégories de tumeurs cérébrales.

\begin{table}[htbp]
  \centering
  \caption{Répartition des images par classe.}
  \label{tab:dataset-classes}
  \begin{tabular}{@{}lc@{}}
    \toprule
    \thd{Classe} & \thd{Nombre d'images} \\
    \midrule
    Glioma          & 500 \\
    Meningioma      & 500 \\
    Pituitary Tumor & 500 \\
    \midrule
    \textbf{Total}  & \textbf{\NTotal} \\
    \bottomrule
  \end{tabular}
\end{table}

Cette répartition équilibrée permet d'assurer une évaluation plus fiable des performances
des modèles de classification en évitant qu'une classe ne domine le processus
d'apprentissage.

\subsection{Description des classes}

Les trois catégories étudiées correspondent aux types de tumeurs cérébrales les plus
fréquemment rencontrés dans les jeux de données publics d'IRM~:

\begin{itemize}[leftmargin=1.4em,itemsep=2pt]
  \item \textbf{Glioma}~: tumeur développée à partir des cellules gliales du cerveau. Elle
        représente l'une des formes les plus courantes de tumeurs cérébrales primaires.
  \item \textbf{Meningioma}~: tumeur provenant des méninges, les membranes entourant le
        cerveau et la moelle épinière. Elle est généralement bénigne mais peut entraîner des
        complications selon sa localisation.
  \item \textbf{Pituitary Tumor}~: tumeur affectant l'hypophyse, une glande endocrine
        située à la base du cerveau. Son diagnostic précoce est essentiel afin de prévenir
        les perturbations hormonales associées.
\end{itemize}

L'objectif du projet consiste à identifier automatiquement la catégorie à laquelle
appartient chaque image IRM à partir des caractéristiques visuelles extraites.

\begin{figure}[htbp]
  \centering
  \includegraphics[width=\textwidth]{dataset}
  \caption{Vue d'ensemble du dataset~: répartition train/test par classe et échantillons
    visuels des trois catégories de tumeurs.}
  \label{fig:dataset-overview}
\end{figure}

% -----------------------------------------------------------------------
\section{Répartition des données d'apprentissage et de test}

Afin d'évaluer correctement la capacité de généralisation des modèles, le jeu de données a
été divisé en deux sous-ensembles distincts~: un ensemble d'apprentissage destiné à
l'entraînement des modèles et un ensemble de test réservé à l'évaluation finale.

Une stratégie de division stratifiée (\emph{Stratified Train-Test Split}) a été adoptée afin
de préserver la proportion de chaque classe dans les deux sous-ensembles.

La répartition retenue est de~:

\begin{itemize}[leftmargin=1.4em,itemsep=2pt]
  \item \textbf{80\,\% pour l'apprentissage}~;
  \item \textbf{20\,\% pour le test}.
\end{itemize}

Cette configuration conduit à la distribution présentée dans le Tableau~\ref{tab:train-test}.

\begin{table}[htbp]
  \centering
  \caption{Répartition Train/Test.}
  \label{tab:train-test}
  \begin{tabular}{@{}lcc@{}}
    \toprule
    \thd{Ensemble} & \thd{Nombre d'images} & \thd{Pourcentage} \\
    \midrule
    Train & \NTrain & 80\,\% \\
    Test  & \NTest  & 20\,\% \\
    \midrule
    Total & \NTotal & 100\,\% \\
    \bottomrule
  \end{tabular}
\end{table}

Grâce à cette stratégie, chaque classe reste représentée de manière équilibrée dans les
ensembles d'apprentissage et de test, garantissant ainsi une évaluation impartiale des
performances.

% -----------------------------------------------------------------------
\section{Prétraitement des images}

Avant l'extraction des caractéristiques, plusieurs opérations de prétraitement ont été
appliquées afin d'améliorer la qualité des images et d'homogénéiser les données. Ces étapes
visent à réduire les variations non pertinentes et à faciliter l'extraction de descripteurs
robustes.

\subsection{Conversion en niveaux de gris}

Les images IRM utilisées dans cette étude sont converties en niveaux de gris afin de
simplifier leur représentation tout en conservant les informations anatomiques essentielles.

Cette transformation permet de réduire la dimension des données et d'uniformiser les
traitements ultérieurs. Elle est particulièrement adaptée aux images médicales IRM où les
informations diagnostiques sont principalement contenues dans les variations d'intensité
plutôt que dans les composantes colorimétriques.

\subsection{Redimensionnement des images}

Afin de garantir une représentation uniforme des données, toutes les images sont
redimensionnées vers une résolution fixe avant l'extraction des caractéristiques.

Cette étape permet~:

\begin{itemize}[leftmargin=1.4em,itemsep=2pt]
  \item d'assurer la cohérence des dimensions entre les échantillons~;
  \item de faciliter l'extraction des descripteurs~;
  \item de réduire le coût computationnel des traitements.
\end{itemize}

Le redimensionnement garantit également que les caractéristiques extraites soient
comparables entre toutes les images du dataset.

\subsection{Normalisation des intensités}

Les images IRM peuvent présenter des variations importantes d'intensité liées aux conditions
d'acquisition ou aux différences entre appareils médicaux.

Une étape de normalisation est donc appliquée afin de ramener les niveaux de gris dans une
plage commune de valeurs. Cette opération contribue à réduire l'influence des variations
d'éclairage et améliore la robustesse des caractéristiques extraites.

La normalisation facilite également l'apprentissage des modèles en garantissant une échelle
homogène des données.

\subsection{Amélioration du contraste par CLAHE}

Afin de renforcer la visibilité des structures anatomiques et des régions tumorales, une
méthode d'amélioration locale du contraste appelée \textbf{Contrast Limited Adaptive
Histogram Equalization} (CLAHE) est appliquée.

Contrairement à l'égalisation d'histogramme classique, CLAHE agit localement sur différentes
régions de l'image tout en limitant l'amplification du bruit.

Cette technique présente plusieurs avantages~:

\begin{itemize}[leftmargin=1.4em,itemsep=2pt]
  \item amélioration du contraste local~;
  \item meilleure mise en évidence des textures~;
  \item renforcement des contours et structures tumorales~;
  \item amélioration de la qualité des caractéristiques extraites.
\end{itemize}

L'application de CLAHE s'avère particulièrement bénéfique pour les descripteurs de texture
tels que GLCM et LBP, qui reposent fortement sur les variations locales d'intensité.

\subsection{Pipeline complet de prétraitement}

L'ensemble des opérations réalisées avant l'extraction des caractéristiques est résumé dans
la Figure~\ref{fig:preprocessing-pipeline}.

\begin{notebox}[Pipeline de prétraitement~:]
Image IRM originale $\rightarrow$ Conversion en niveaux de gris $\rightarrow$
Redimensionnement $\rightarrow$ Normalisation des intensités $\rightarrow$
Amélioration du contraste (CLAHE) $\rightarrow$ Extraction des caractéristiques
(Statistiques, GLCM, LBP, DWT, HOG).
\end{notebox}

\begin{figure}[htbp]
  \centering
  \includegraphics[width=0.92\textwidth]{preprocessing_pipeline}
  \caption{Pipeline complet de prétraitement des images IRM avant extraction des
    caractéristiques.}
  \label{fig:preprocessing-pipeline}
\end{figure}

Ce pipeline permet d'obtenir des images homogènes et de meilleure qualité, favorisant ainsi
l'extraction de caractéristiques discriminantes et l'amélioration des performances de
classification.
